% !TeX TS-program = pdflatex
\documentclass[14pt, a4paper, oneside]{extreport}
\usepackage{extsizes} % Поддержка 14pt

% Кодировки и язык
\usepackage[T2A]{fontenc}
\usepackage[utf8]{inputenc}
\usepackage[russian]{babel}

% Геометрия страницы (ГОСТ Р 7.0.11-2011)
\usepackage{geometry}
\geometry{left=30mm, right=10mm, top=20mm, bottom=20mm}

% Интервалы и абзацы
\usepackage{setspace}
\onehalfspacing % 1.5 межстрочный интервал
\setlength{\parindent}{1.25cm}
\setlength{\parskip}{0pt}
\usepackage{indentfirst}

% Математика и формулы
\usepackage{amsmath, amssymb, mathtools}
\numberwithin{equation}{section} % Нумерация формул: (1.1), (2.3) и т.д.

% Рисунки, таблицы, подписи
\usepackage{graphicx}
\usepackage{caption}
\captionsetup{labelsep=period, font=normal}

% Заголовки и нумерация разделов
\usepackage{titlesec}
\titleformat{\chapter}[hang]{\bfseries\centering\large}{\thechapter}{1em}{}
\titleformat{\section}[hang]{\bfseries\normalsize}{\thesection}{1em}{}
\titleformat{\subsection}[hang]{\bfseries\normalsize}{\thesubsection}{1em}{}
\titlespacing*{\chapter}{0pt}{0pt}{1em}
\titlespacing*{\section}{0pt}{1em}{0.5em}
\titlespacing*{\subsection}{0pt}{1em}{0.5em}

% Убираем слово "Глава", оставляем только номер
\addto\captionsrussian{\renewcommand{\chaptername}{}}
\renewcommand{\contentsname}{СОДЕРЖАНИЕ}

% Ссылки
\usepackage{hyperref}
\hypersetup{colorlinks=false, pdfborder={0 0 0}}

\begin{document}

% ==========================================
% НУМЕРАЦИЯ НАЧИНАЕТСЯ С 6-Й СТРАНИЦЫ
% ==========================================
\pagenumbering{arabic}
\setcounter{page}{6}

\tableofcontents
\clearpage

% ==========================================
% ВВЕДЕНИЕ
% ==========================================
\chapter*{ВВЕДЕНИЕ}
\addcontentsline{toc}{chapter}{ВВЕДЕНИЕ}
% ВСТАВЬТЕ СЮДА ТЕКСТ ВВЕДЕНИЯ
\clearpage

% ==========================================
% ГЛАВА 1
% ==========================================
\chapter{АНАЛИЗ ПРЕДМЕТНОЙ ОБЛАСТИ И СУЩЕСТВУЮЩИХ АНАЛОГОВ}
\section{Обзор исследований и моделей машинного обучения для решения аналогичных задач}
\subsection{Применение классических и глубоких автокодировщиков для оценки состояния оборудования}
% Текст 1.1.1
\subsection{Вариационные автокодировщики в задачах вероятностного прогнозирования и генерации временных рядов}
% Текст 1.1.2
\subsection{Сравнительный анализ генеративных архитектур (GAN, Diffusion, VAE) для моделирования циклических процессов}
% Текст 1.1.3
\subsection{Адаптация функции ELBO и амортизированный вывод для многошагового прогнозирования состояний}
% Текст 1.1.4

\section{Критерии оценки и валидации генеративных моделей в промышленной диагностике}
% Текст 1.2

\subsection*{Научная новизна}
\addcontentsline{toc}{subsection}{Научная новизна}
\begin{itemize}
    \item Предложена адаптация вариационной нижней оценки...
    \item Разработан метод формирования скалярного индикатора...
    \item Создан двухэтапный алгоритм генерации...
\end{itemize}

\subsection*{Теоретическая и практическая значимость}
\addcontentsline{toc}{subsection}{Теоретическая и практическая значимость}
% Текст значимости

\subsection*{Положения, выносимые на защиту}
\addcontentsline{toc}{subsection}{Положения, выносимые на защиту}
\begin{enumerate}
    \item Математическая модель оценки состояния...
    \item Алгоритм двухэтапной генерации...
    \item Методика валидации генеративной модели...
\end{enumerate}

\section*{1.3 Выводы}
\addcontentsline{toc}{section}{1.3 Выводы}
% Текст выводов
\clearpage

% ==========================================
% ГЛАВА 2
% ==========================================
\chapter{ПРОЕКТИРОВАНИЕ СИСТЕМЫ МОДЕЛИРОВАНИЯ ДАННЫХ}
\section{Постановка задачи}
\subsection{Формальная постановка задачи моделирования состояния оборудования}
% Текст 2.1.1 + формулы
\begin{equation}
X_{1:n} = \{\mathbf{x}_1, \mathbf{x}_2, \dots, \mathbf{x}_n\}
\end{equation}

% Добавьте остальные подразделы главы 2 по необходимости
\clearpage

% ==========================================
% ЗАКЛЮЧЕНИЕ
% ==========================================
\chapter*{ЗАКЛЮЧЕНИЕ}
\addcontentsline{toc}{chapter}{ЗАКЛЮЧЕНИЕ}
% Текст заключения
\clearpage

% ==========================================
% СПИСОК ИСТОЧНИКОВ
% ==========================================
\begin{thebibliography}{99}
\addcontentsline{toc}{chapter}{СПИСОК ИСПОЛЬЗОВАННЫХ ИСТОЧНИКОВ}
\bibitem{1} Пример источника...
\bibitem{2} Jia Guo, G. Liu, Y. Zuo, J. Wu. An Anomaly Detection Framework...
% Добавьте остальные источники
\end{thebibliography}

\end{document}